\documentclass[10pt,twocolumn,letterpaper]{article}

\usepackage[letterpaper,margin=0.85in]{geometry}
\usepackage{titling}
\setlength{\columnsep}{0.3in}
\setlength{\droptitle}{-3em}
\renewcommand{\abstractname}{\large Abstract}
\usepackage{cite}
\usepackage{amsmath,amssymb,amsfonts}
\usepackage{graphicx}
\usepackage{textcomp}
\usepackage{xcolor}
\usepackage{booktabs}
\usepackage{multirow}
\usepackage{url}
\usepackage[breaklinks=true,hidelinks]{hyperref}
\usepackage{caption}
\captionsetup{font=footnotesize}

\title{Trade Market Manipulation Detection with\\
Graph Neural Networks:\\
A Calibrated-Synthesis-to-Deployment\\
Surveillance Pipeline}

\author{Sumit Sontakke\\
\small Department of Computer Science, PES University, Bengaluru, India\\
\small \texttt{pes2pge24ds111@stu.pes.edu}}

\begin{document}
\maketitle

\begin{abstract}
We describe an end-to-end surveillance pipeline for detecting collusive
trade-market manipulation that ties together five components: a daily
NSE bhavcopy calibration loop, a Cont-2001-stylised synthetic market
generator, an agent-based simulator (ABIDES) used as a second substrate
to test generator-shift, a graph neural network detector with an
optional second-stage classifier, and a Streamlit investigation
dashboard that surfaces flagged traders with per-feature evidence and a
language-model justification card. We argue that the pipeline is the
contribution: each layer was driven by a specific surveillance
requirement (interpretability for investigators, robustness across
parameter shifts, controllable false-positive cost), and the layers
interlock through a small set of explicit interfaces. We instantiate
the detector layer with three architectures---a 2-feature GraphSAGE
baseline, a gradient-boosted bolt-on classifier consuming the GraphSAGE
trader score plus six engineered features, and a feature-augmented
GraphSAGE that ingests the same six features as node attributes---and
report results across within-generator, cross-generator, and
leave-one-family-out regimes. The augmented network lifts within-OOD
edge AUC from \(0.638\) to \(0.984\); the bolt-on classifier reaches
pooled trader AUC \(0.968\); the augmented network retains AUC
\(0.842\) across a different simulator while the baseline collapses to
chance. We discuss the practical implication: per-generator threshold
calibration, an investigator dashboard, and an interpretable feature
representation are at least as important as model complexity.
\end{abstract}

\vspace{0.5em}\noindent\textbf{Keywords:}
Market surveillance, collusive trading detection, graph neural networks,
calibrated synthesis, agent-based simulation, investigation dashboard,
deployment, operating-point selection.

\section{Introduction}
\label{sec:intro}

Trade-market surveillance systems sit at the intersection of three
distinct engineering pressures. They must satisfy a regulatory consumer
that wants explicit, defensible rules; they must satisfy an investigator
team whose workload determines the operationally acceptable purity
floor; and they must satisfy a machine-learning team that wants to
exploit modern representational learning to lift recall beyond what
hand-crafted rules can reach. Each of these constituents pulls in a
different direction. The investigator team wants per-flag evidence and
a clean operating point at high precision. The regulator wants the rule
or model to be auditable and stable. The ML team wants the freedom to
use whatever architecture lifts AUC the most.

This paper does not present a single best model. Instead it presents a
pipeline---five interlocking components---designed so that the three
pressures coexist. The pipeline begins with daily NSE bhavcopy
ingestion and calibration of a synthetic market generator to the
empirical first moments of the index. The generator is then driven
with three manipulation taxonomies (clique, ring, mixed) at controllable
parameter envelopes to produce labelled training cohorts that are both
realistic enough to transfer and varied enough to expose distribution
shift. A second substrate (the ABIDES agent-based simulator) is used in
parallel as a cross-generator stress test. A graph neural network
backbone, optionally augmented with a second-stage classifier, sits at
the centre of the detection layer; six explicitly engineered per-trader
features provide both the lift and the interpretability hooks the
investigator team needs. A Streamlit dashboard surfaces the flagged
traders, the price/volume context of the manipulated ticker, the
counterparty network, and a per-trader language-model justification.

We claim three contributions. First, we describe the pipeline as a
whole, including the data interfaces and the small set of design
decisions that connect each layer. Second, we report empirical results
across three increasingly demanding cohorts, framed as system tests:
within-generator out-of-distribution, cross-generator (a different
simulator entirely), and a leave-one-family-out variant in which one
manipulation taxonomy is held out from training. Third, we describe the
operating-point and threshold-calibration logic that we believe is the
correct deployment hygiene for a model whose ranking is stable across
generators but whose absolute score distribution is not.

The paper proceeds as follows. Section~\ref{sec:related} situates the
work against existing market-surveillance and graph-fraud literature.
Section~\ref{sec:overview} describes the pipeline at the system level
and lists the interfaces between components. Section~\ref{sec:data}
describes the data pipeline (calibration, synthesis, ABIDES
integration). Section~\ref{sec:features} describes the six engineered
features as surveillance signals. Section~\ref{sec:detector} describes
the three detector variants we compare. Section~\ref{sec:projection}
discusses the trader-projection and operating-point logic.
Section~\ref{sec:dashboard} sketches the investigator dashboard.
Section~\ref{sec:results} reports empirical validation across the three
test cohorts. Section~\ref{sec:deploy} discusses deployment
considerations. Section~\ref{sec:lim} covers limitations and
Section~\ref{sec:concl} concludes.

\section{Related Work}
\label{sec:related}

The literature on market-manipulation detection has historically been
dominated by rule-driven systems that emit alerts when explicit
microstructure conditions hold: layering thresholds on order
cancellation rates, spoofing tests on persistence of unfilled top-book
orders, pump-and-dump signatures on price-volume spikes. The rule era
delivered high false-positive rates; the shift to supervised classifiers
came later, motivated by the
observation that rule-based systems have constant false-positive cost
per rule and therefore scale poorly when many rules are layered.
Cao et al.\ \cite{cao2014pso} applied particle-swarm-optimised
classifiers to labelled synthetic and exchange data, treating
manipulation detection as a per-order or per-trader binary
classification.

The graph turn is more recent. The genuinely collusive activity that
modern markets observe---ring trading, clique cross-trading, front
account schemes---does not have a per-trader signature. Each individual
account in a ring may look unremarkable; the conspiracy lives in the
edges. Graph neural networks have therefore become a natural fit for
graph-based financial fraud detection. We use GraphSAGE
\cite{hamilton2017graphsage} for its inductive properties: production
deployments need to operate on graphs whose node identities differ
between training and test sessions. Spectral approaches such as GCN
\cite{kipf2017gcn} are widely used but less natural in the inductive
setting, and attention-based architectures such as GAT
\cite{velickovic2018gat} could in principle be substituted for the
backbone we adopt.

Cross-distribution evaluation has been emphasised in image
classification by Recht et al.\ \cite{recht2019imagenet} and made
systematic in the WILDS \cite{koh2021wilds} benchmark. The analogue in
surveillance is to vary not only the seeds and parameter envelope but
also the underlying simulator; we use ABIDES \cite{byrd2019abides} as
a structurally different second generator for that test. For the
synthesiser's stylised facts we follow Cont \cite{cont2001stylized}.

\section{Pipeline Overview}
\label{sec:overview}

The pipeline comprises five layers shown schematically below. Each
layer communicates with the next through an explicit, versioned
artifact, which keeps each layer independently testable and
re-runnable.

\smallskip
\noindent
\fbox{\parbox{0.94\linewidth}{\footnotesize
\begin{enumerate}
\item Daily NSE bhavcopy fetch and calibration $\rightarrow$
  parameters of stylised-facts synthesiser
\item Synthesiser run (with optional ABIDES substrate)
  $\rightarrow$ \texttt{orders.csv} and \texttt{trades.csv} per scenario
\item Graph construction and feature extraction
  $\rightarrow$ per-run PyG graphs and per-trader feature CSVs
\item Detector inference (baseline / bolt-on / augmented)
  $\rightarrow$ per-trader probability scores
\item Investigation dashboard (Streamlit + LLM)
  $\rightarrow$ investigator-facing flag list with evidence
\end{enumerate}}}
\smallskip

Each artifact is small, versioned by cohort, and cached. Re-running a
single layer is cheap: feature extraction over a 50-run cohort
completes in under two minutes on CPU; the bolt-on classifier under
leave-one-run-out cross-validation completes in under \(45\) seconds;
inference on a 50-run cohort takes less than a second on GPU. This
makes the pipeline iterable enough that ablations and operating-point
sweeps are routine rather than expensive.

\section{Data Pipeline}
\label{sec:data}

\subsection{NSE Calibration Loop}
The data pipeline begins with daily NSE bhavcopy ingestion. A small
fetcher service downloads end-of-day bhavcopy files for the targeted
ticker universe and stores them in SQLite. A calibrator service
estimates the first and second moments of log-returns over a rolling
window and emits a calibration manifest that the synthesiser ingests as
its base regime. This manifest is the only direct link between real
data and synthetic data; everything downstream is generated, labelled
and varied within the controllable envelope.

\subsection{Primary Synthesiser}
The primary synthesiser produces a daily session of orders and trades
for a configurable population (typically 500 traders). It follows the
Cont stylised facts: heavy-tailed log-returns, volatility clustering,
absence of linear return autocorrelation. Manipulation scenarios are
injected on top of an otherwise benign session via three families:
\emph{clique} (a small set of accounts pushing or absorbing price
through a shared front), \emph{ring} (a closed cycle of accounts
passing positions around), and \emph{mixed} (a session combining both
patterns). The injection layer is parameterised by core manipulator
count, manipulation ratio, duration, and concealment level---giving us
the controllable envelope we use to define training, evaluation, and
cross-distribution cohorts.

\subsection{ABIDES Substrate}
A second substrate is the ABIDES agent-based discrete-event simulator
\cite{byrd2019abides}, used here as a cross-generator stress test
rather than a primary training source. We implemented three ABIDES
agents (\texttt{CollusiveCliqueAgent}, \texttt{RingTraderAgent},
\texttt{FrontAccountAgent}) that approximate the same three taxonomies
used in the primary synthesiser, but whose underlying microstructure
(limit order book dynamics, exchange agent behaviour, message
ordering) differs substantially. ABIDES sessions are then run through
the same downstream pipeline as primary-synthesiser sessions; an
adapter writes ABIDES output to the standard
\texttt{orders.csv}/\texttt{trades.csv}/\texttt{scenarios.csv}/
\texttt{traders.csv} schema so that the detector and dashboard layers
are generator-agnostic.

\subsection{Cohort Definitions}
Three cohorts are used in this paper. The training cohort consists of
five days at \(\sim\!50\) runs per day, drawn from a wide parameter
envelope. The within-generator OOD cohort consists of \(50\) runs from
the primary synthesiser, drawn from parameter ranges strictly outside
the training envelope (session duration \(60\)--\(90\) min,
manipulation ratio \(2\)--\(5\%\), core count \(16\)--\(25\)) with
seeds disjoint from training. The cross-generator cohort consists of
\(15\) ABIDES runs (\(5\) seeds across \(3\) families on a held-out
calibration date), supplementing a \(3\)-run pilot. Table
\ref{tab:cohorts} summarises.

\begin{table}[t]
\caption{Cohort summary.}
\label{tab:cohorts}
\centering
\footnotesize
\begin{tabular}{lccc}
\toprule
Cohort & Runs & Traders/run & Manip./run \\
\midrule
Training (5-day, primary) & 250 & $\sim$500 & $\sim$10--20 \\
Within-generator OOD & 50 & $\sim$500 & 16--25 \\
ABIDES pilot & 3 & $\sim$125 & 6 \\
ABIDES expanded & 15 & $\sim$125 & 6 \\
\bottomrule
\end{tabular}
\end{table}

\section{Feature Engineering for Surveillance}
\label{sec:features}

The pipeline includes six per-trader features computed over each
trader's busiest five-minute order window. Each feature was chosen so
that an investigator can be told, in one sentence, what manipulation
behaviour it captures. This is deliberate: a feature whose meaning a
human cannot articulate is not deployable in a surveillance setting,
even if it is statistically powerful. Below, \(O_v\) denotes the order
set of trader \(v\); \(B_v\) is the five-minute window with the
maximum order count within \(O_v\); \(T_v\) is the set of trades in
which \(v\) participated during \(B_v\); \(s_k\) is the share of \(T_v\)
attributable to counterparty \(k\); and \(q\) is the vector of order
quantities placed by \(v\).

\begin{enumerate}
\item \textbf{Burst concentration}, \(\phi_1(v) = |B_v|/|O_v|\). A
ring or clique that spikes activity in a tight window produces
\(\phi_1 \gg 0\); a steadily active legitimate trader produces
\(\phi_1 \approx 0.1\).
\item \textbf{Side entropy in burst}, \(\phi_2(v) =
-p\log_2 p - (1{-}p)\log_2(1{-}p)\), where \(p\) is the buy fraction
inside \(B_v\). A pump tends to be one-sided ($\phi_2 \to 0$); a
ring rotating positions tends to be balanced ($\phi_2 \to 1$).
\item \textbf{Counterparty HHI}, \(\phi_3(v) = \sum_k s_k^2\). A
front-account scheme concentrates burst trades against one or two
favoured counterparties; ring members deliberately spread.
\item \textbf{Order quantity coefficient of variation},
\(\phi_4(v) = \sigma(q)/\mu(q)\). Synthetic manipulators often place
uniformly-sized orders; benign traders vary.
\item \textbf{Top-partner trade share}, \(\phi_5(v) = \max_k s_k\). A
direct readout of the strongest single counterparty link.
\item \textbf{Co-active top-quantile count}, \(\phi_6(v)\): the number
of other top-activity-quantile traders whose burst windows overlap that
of \(v\). A structural marker for coordinated activity.
\end{enumerate}

The standalone power of each feature on the primary OOD cohort is
summarised in Table~\ref{tab:standalone}. Importantly, the
\emph{direction} of each feature is learned from data rather than
hand-coded: on this cohort, the manipulator direction for
\(\phi_3\) is in fact \emph{low} (because the ring scenarios spread
counterparties), which is the opposite of the textbook intuition. The
ensemble downstream therefore needs to be sign-agnostic; this is the
specific reason we use a gradient-boosted tree ensemble as the bolt-on
classifier rather than a linear model.

\begin{table}[t]
\caption{Standalone ROC AUC of engineered features (sign-agnostic).}
\label{tab:standalone}
\centering
\footnotesize
\begin{tabular}{lcc}
\toprule
Feature & Standalone AUC & Direction \\
\midrule
$\phi_3$ Counterparty HHI in burst & 0.800 & $-$ \\
$\phi_5$ Top-partner trade share & 0.775 & $-$ \\
$\phi_6$ Co-active top-quantile count & 0.744 & $-$ \\
$\phi_1$ Burst concentration & 0.704 & $-$ \\
$\phi_4$ Order qty coefficient of variation & 0.666 & $+$ \\
$\phi_2$ Side entropy in burst & 0.550 & $-$ \\
\bottomrule
\end{tabular}
\end{table}

\section{Detector Stack}
\label{sec:detector}

The pipeline's detector layer is deliberately modular. Three variants
are evaluated.

\subsection{Baseline GraphSAGE (v1)}
A two-layer GraphSAGE network with \texttt{max} aggregator, hidden
dimensions \((256,128,64)\), dropout \(0.3\). The input is the directed
trade graph of one session; the output is a per-edge manipulation
probability. Node features are two-dimensional---trader's total
executed quantity and number of distinct counterparties---which we
treat as a graph-topology baseline. Edge features are
eight-dimensional, covering lead-lag correlation, signed imbalance,
traded volume, interaction count, pre- and post-window mid-price drift,
and VWAP imbalance. The model trains with focal loss
(\(\alpha{=}0.85,\gamma{=}2.0\)) and a five-day continual warm-start
schedule. Early stopping is governed by per-family validation trader
recall rather than validation cross-entropy, in keeping with the
surveillance operating metric.

\subsection{Tier-2 Bolt-On Classifier}
The second variant attaches a gradient-boosted classifier to the
baseline's per-trader score. The classifier consumes a seven-element
input vector: the baseline GraphSAGE's trader score plus the six
engineered features of \S\ref{sec:features}. Hyperparameters are
\(n=120\) estimators, depth \(3\), learning rate \(0.08\), subsample
\(0.85\). The classifier is trained under leave-one-run-out
cross-validation over the \(50\)-run OOD cohort, producing pooled
per-trader probabilities that we then sweep for an operating point.
This bolt-on stack is the simplest, cheapest, and most interpretable
configuration; it inherits the baseline GraphSAGE checkpoint without
modification and adds 40 seconds of CPU compute.

\subsection{Feature-Augmented GraphSAGE (v4)}
The third variant feeds the six engineered features into the network
itself as additional node attributes. The node-feature vector grows
from two dimensions to eight, and the network is retrained from
scratch (the first convolution layer's weight matrix changes shape,
precluding warm-start from the baseline checkpoint). All other
hyperparameters and the training schedule are held fixed, so the
only variable between the baseline and the augmented network is the
node-feature set.


\subsection{Implementation Details}
\label{sec:impl}
All graph network training and inference runs in a Docker container
provisioned with PyTorch~2.x and PyTorch Geometric, on an NVIDIA RTX
\(5080\) workstation card (sm\_120, \(16\) GB VRAM). The bolt-on
classifier uses scikit-learn's \texttt{GradientBoostingClassifier}.
Graph construction and feature extraction are pure NumPy/pandas and
are CPU-portable. Hyperparameters held constant across all comparisons
are listed in Table~\ref{tab:hparams}. A single five-day continual
training run for the augmented network takes approximately
\(170\)~minutes on the workstation card; feature caching for the OOD
cohort completes in under two minutes on CPU; the bolt-on classifier
under leave-one-run-out cross-validation (\(50\) folds, $\sim$10{,}000
trader-rows) completes in under \(45\) seconds on a single CPU core;
inference on the \(50\)-run OOD cohort for either network takes well
under one second on GPU.

\begin{table}[t]
\caption{Hyperparameters held constant across all comparisons.}
\label{tab:hparams}
\centering
\footnotesize
\begin{tabular}{lll}
\toprule
Component & Parameter & Value \\
\midrule
\multirow{6}{*}{GraphSAGE} & layers & 2 \\
 & hidden dims & 256, 128, 64 \\
 & aggregator & max \\
 & dropout & 0.3 \\
 & loss & focal $(\alpha{=}0.85,\gamma{=}2.0)$ \\
 & early-stop patience & 6 (val trader recall) \\
\midrule
\multirow{4}{*}{Tier-2 GBM} & n\_estimators & 120 \\
 & max\_depth & 3 \\
 & learning\_rate & 0.08 \\
 & subsample & 0.85 \\
\midrule
Both & random seed & 42 \\
\bottomrule
\end{tabular}
\end{table}

\subsection{Pipeline Orchestration}
\label{sec:orch}
The five pipeline layers communicate through file-based artifacts
rather than in-memory handoff, which keeps each layer independently
runnable and makes the pipeline easy to resume after a crash.
The training loop is summarised below.

\begin{figure}[t]
\small
\noindent\textbf{Pipeline orchestration (one cohort end-to-end):}
\begin{enumerate}
\item Fetch NSE bhavcopy for the trading day and persist to SQLite.
\item Calibrate the synthesiser parameters from the latest bhavcopy
window; persist a calibration manifest.
\item Generate the cohort using the calibrated synthesiser (or ABIDES);
write \texttt{orders.csv}, \texttt{trades.csv}, \texttt{scenarios.csv},
\texttt{traders.csv} per run.
\item Extract per-run engineered features and cache to
\texttt{outputs/\_phase\_g\_features/}.
\item Build per-run graphs (cached to disk after first build).
\item For each detector variant: load matching graphs, train or load
checkpoint, score per trader, persist to
\texttt{outputs/\_phase\_g\_v\{N\}\_trader\_scores/}.
\item Sweep operating-point thresholds; persist a per-detector
operating-points JSON for the dashboard.
\item Launch the Streamlit dashboard pointing at the cohort directory.
\end{enumerate}
\end{figure}

The atomic step in the loop is step 6: every other step caches its
output. A single ablation (re-running step 6 with a different detector
variant) typically completes in tens of seconds on cached features and
graphs.
\section{Trader Projection and Operating-Point Selection}
\label{sec:projection}

The detector layer emits per-edge probabilities; the dashboard and the
investigator workload calculus need per-trader scores. We use a
weighted projection of each trader's incident-edge probabilities,
\begin{equation}
\hat{y}(v) = 0.7\,\max_{e\in E(v)} p_e
\;+\;0.3\,\frac{1}{3}\!\!\sum_{e\in\mathrm{top}_3(v)}\!\!p_e ,
\label{eq:proj}
\end{equation}
which balances the strong-signal bias of a pure \texttt{max} (which
over-weights single noisy edges and hurts purity) against the
dilution of a pure mean.

The threshold applied to \(\hat{y}\) is then chosen by an
investigator-workload calculus rather than by a single statistical
criterion. For an OOD cohort whose pooled labels are known, sweeping
the threshold produces a recall/purity curve with several natural
operating points. We surface three of them as defaults
(Table~\ref{tab:ops}): a high-precision point useful when investigator
capacity is constrained, a balanced point at the best F\(_1\), and a
high-recall point useful when investigator capacity is loose. The
dashboard exposes a slider so that an investigator can move along the
curve and watch the per-flag count change in real time.

\begin{table}[t]
\caption{Operating points for the tier-2 classifier on the primary OOD cohort.}
\label{tab:ops}
\centering
\footnotesize
\begin{tabular}{lccc}
\toprule
Operating point & Threshold & Recall & Purity \\
\midrule
High-precision & 0.86 & 0.619 & 0.991 \\
Best F$_1$ & 0.48 & 0.759 & 0.941 \\
High-recall & 0.18 & 0.850 & 0.803 \\
\bottomrule
\end{tabular}
\end{table}

\section{Investigation Dashboard}
\label{sec:dashboard}

The dashboard is built on Streamlit and consumes only the artifacts
emitted by the layers above. Its job is not to retrain or score---all
scoring is pre-computed---but to assemble each flagged trader's
evidence in a single view. Six sections appear in sequence:

\begin{enumerate}
\item \textbf{Run selector and cohort context}. A dropdown of the
fifty OOD runs with a one-line summary (family, seed, parameter slice).
\item \textbf{Feature library}. The six features with their
descriptions, the live sample value for a flagged trader, and the
standalone AUC.
\item \textbf{Flagged-trader 2D scatter}. Trader score on one axis, a
selectable feature on the other; the threshold slider drags a dashed
horizontal line and the flagged set updates instantly.
\item \textbf{Manipulated-ticker price and volume}. One-minute mid
price and volume of the ticker the manipulator operated on, with a
shaded band marking the burst window.
\item \textbf{Counterparty network}. A two-ring graph showing the
top-flagged traders inside and their counterparties outside, with edge
thickness proportional to trade count.
\item \textbf{LLM justification}. For each flagged trader, a short
sentence generated by a local model (qwen2.5:7b) constrained by the
feature values, the per-trader trade count, and the ground-truth label
where available. The model paraphrases the numbers into prose; it does
not introduce evidence that the numbers do not support.
\end{enumerate}

In practice the dashboard is consulted both at training time (to
identify which features are doing what work, often by visual
inspection of the scatter) and at deployment time (to give an
investigator a single page per flagged trader with all relevant
evidence laid out).


\subsection{Dashboard Layout}
\label{sec:dashlayout}
Figure~\ref{fig:dash} sketches the dashboard layout. The flow is
top-to-bottom: cohort selection $\rightarrow$ feature library
$\rightarrow$ flagged-trader scatter $\rightarrow$ price/volume context
$\rightarrow$ counterparty network $\rightarrow$ per-trader LLM
justification. Each chart is accompanied by a small explanatory box,
generated once per cohort and cached for one hour, that paraphrases
the chart's specific numbers in plain English. This is a deliberate
design choice: investigators tend to skim charts, and a one-sentence
plain-language summary substantially improves the signal that the
chart's contents are absorbed before the investigator moves on.

\begin{figure}[t]
\small
\noindent\fbox{\parbox{0.94\linewidth}{
\textbf{Streamlit dashboard:} \\
{\small
A. Run selector + cohort metrics\\
B. Feature library (6 features, sample value, AUC) \\
C. Flagged-trader 2D scatter + threshold slider\\
D. Manipulated-ticker price/volume with burst-window band\\
E. Counterparty network (focal traders + counterparties)\\
F. Per-trader LLM justification cards \\
}
}}
\caption{Dashboard layout (per-cohort view). Each section consumes only
pre-computed artifacts; no model inference happens inside Streamlit.}
\label{fig:dash}
\end{figure}

\subsection{LLM Justification Layer}
\label{sec:llm}
The LLM justification card is generated by a local model
(\texttt{qwen2.5:7b} via Ollama) prompted with the flagged trader's
six feature values, total order count, in-burst order count, and the
ground-truth label where available. The prompt is structured so that
the model paraphrases the supplied facts into a two-sentence verdict;
it is not given a free hand to introduce evidence the numbers do not
support. Cards are cached for one hour, keyed on the prompt string,
so the dashboard pays the model cost once per cohort even if multiple
investigators consult it. The card is intended as a \emph{speech-act
layer}: it converts numbers into prose, lowering the cognitive cost
for an investigator who needs to triage many flags quickly.

\section{Empirical Validation}
\label{sec:results}

We validate the pipeline against three increasingly demanding cohorts.

\subsection{Within-Generator OOD}
Table~\ref{tab:ood} summarises the primary result. The baseline
GraphSAGE achieves edge AUC \(0.638\), per-family trader recall
\(0.41\)--\(0.48\), and per-family purity \(0.09\)--\(0.14\)---a
clearly inadequate operating point. The augmented network achieves
edge AUC \(0.984\), per-family recall \(0.74\)--\(0.93\), and purity
\(0.71\)--\(0.77\). The bolt-on classifier achieves pooled trader AUC
\(0.968\) and operating points of \(62\%\) recall at \(99\%\) purity
or \(74\%\) recall at \(95\%\) purity.

\begin{table}[t]
\caption{Within-generator OOD results ($n{=}50$ runs).}
\label{tab:ood}
\centering
\footnotesize
\begin{tabular}{lccc}
\toprule
Metric & v1 baseline & Tier-2 GBM & v4 augmented \\
\midrule
Edge AUC & 0.638 & --- & \textbf{0.984} \\
Trader AUC (pooled) & 0.793 & \textbf{0.968} & --- \\
Clique recall & 0.41 & --- & \textbf{0.81} \\
Ring recall & 0.48 & --- & \textbf{0.93} \\
Mixed recall & 0.41 & --- & \textbf{0.75} \\
Recall @ 99\% purity & --- & \textbf{0.62} & --- \\
Recall @ 95\% purity & --- & \textbf{0.74} & --- \\
\bottomrule
\end{tabular}
\end{table}

\subsection{Cross-Generator (ABIDES)}
We re-run inference for the baseline and the augmented network on the
expanded ABIDES cohort (Table~\ref{tab:abides}). The augmented network
retains edge AUC \(0.842\); the baseline collapses to chance
(\(0.518\)). Per-family v4 mean recall remains high (ring
\(1.000\pm0.000\), mixed \(0.983\pm0.015\), clique \(0.905\pm0.113\)).
Purity, however, does not transfer: both detectors land in the
\(0.10\)--\(0.29\) range on ABIDES, much lower than v4's
\(0.71\)--\(0.77\) within-generator. The fact that the baseline and the
augmented network land at the \emph{same} ABIDES purity range argues
that the purity gap is a cohort-level threshold-calibration issue
rather than an architectural weakness of v4. We return to this in
\S\ref{sec:deploy}.

\begin{table}[t]
\caption{Cross-generator (ABIDES expanded, $n{=}15$) result.}
\label{tab:abides}
\centering
\footnotesize
\begin{tabular}{lcc}
\toprule
Metric & v1 baseline & v4 augmented \\
\midrule
Edge AUC & 0.518 (chance) & \textbf{0.842} \\
Clique recall (mean$\pm$sd) & $0.867\pm0.16$ & $\mathbf{0.905\pm0.11}$ \\
Ring recall & $0.864\pm0.03$ & $\mathbf{1.000\pm0.00}$ \\
Mixed recall & $0.837\pm0.07$ & $\mathbf{0.983\pm0.02}$ \\
Clique purity (aggregate) & 0.096 & 0.094 \\
Ring purity & 0.276 & 0.293 \\
Mixed purity & 0.249 & 0.269 \\
\bottomrule
\end{tabular}
\end{table}

\subsection{Leave-One-Family-Out}
The most diagnostic test for the pipeline is to hold out an entire
manipulation family from training and ask the detectors to recognise
it cold. For the bolt-on stack this is a leave-one-family-out
cross-validation over the trader-level data. For the augmented
network it requires three full retrainings, each with one family's
runs removed from the cohort and the per-day validation pool.
Table~\ref{tab:family} reports both.

\begin{table}[t]
\caption{Leave-one-family-out AUC.}
\label{tab:family}
\centering
\footnotesize
\begin{tabular}{lcc}
\toprule
Held-out family & Tier-2 GBM AUC & v4 retrained AUC \\
\midrule
Ring & \textbf{0.994} & 0.480 \\
Clique & 0.975 & 0.533 \\
Mixed & 0.906 & 0.374 \\
\bottomrule
\end{tabular}
\end{table}

The bolt-on classifier holds AUC \(0.91\)--\(0.99\) on the held-out
family. The augmented network collapses to chance or below. We discuss
this contrast in detail in \S\ref{sec:deploy}; the short version is
that the bolt-on stack consumes trader-marginal features, while the
augmented network's message-passing layers learn family-specific
neighbourhood signatures that do not transfer. A methodological caveat
applies to the v4 result: excluding a family from training shrinks the
per-day cohort by a third, which in turn shrinks the validation pool to
three graphs and forces an aggressive early-stop trajectory. A longer
training schedule with a larger validation pool would likely recover
some of the lost AUC, and this remains the most important follow-up
experiment.


\subsection{Feature Ablation Inside the Bolt-On}
\label{sec:abl}
A leave-one-feature-out ablation through the bolt-on classifier
(Table~\ref{tab:abl}) reveals an uneven distribution of load. Dropping
\(\phi_1\) (burst concentration)---only fifth-ranked by standalone
AUC---costs \(-0.051\) AUC and collapses recall at \(99\%\) purity
from \(0.62\) to \(0.27\). Dropping \(\phi_2\) (side entropy), the
weakest feature standalone, costs only \(-0.004\) AUC. The standalone
ranking is therefore a poor proxy for ensemble importance: a feature
that is weak in isolation can be load-bearing in combination, and a
feature that is strong in isolation can be substitutable.

\begin{table}[t]
\caption{Leave-one-feature-out ablation through the bolt-on classifier.}
\label{tab:abl}
\centering
\footnotesize
\begin{tabular}{lcc}
\toprule
Feature dropped & $\Delta$AUC & Recall @ p99 \\
\midrule
\textbf{full (7 features)} & --- & 0.619 \\
$\phi_1$ burst concentration & $-0.051$ & 0.274 \\
$\phi_6$ co-active count & $-0.007$ & 0.592 \\
$\phi_2$ side entropy & $-0.004$ & 0.634 \\
$\phi_4$ order qty CoV & $-0.003$ & 0.608 \\
$\phi_3$ counterparty HHI & $-0.002$ & 0.620 \\
$\phi_5$ top-partner share & $+0.000$ & 0.625 \\
baseline v1 score & $+0.000$ & 0.530 \\
\bottomrule
\end{tabular}
\end{table}

\subsection{Case Study: One Ring Run From OOD}
\label{sec:case}
We pick a representative ring run from the primary OOD cohort and
trace the pipeline's behaviour. The run contains \(\sim\!200\)
traders, six of whom are members of a ring rotating positions among
themselves over a single \(\sim\!45\)-minute window. The six
manipulator traders share three salient properties: burst
concentration \(\phi_1 > 0.85\) (all six, versus a benign median of
\(0.16\)); low counterparty HHI \(\phi_3 < 0.2\), reflecting the
ring's intentionally diversified counterparty list; and high
co-active count \(\phi_6\) because the six rotate trades among
themselves and against each other simultaneously. The dashboard
surfaces all six on the flagged-trader scatter at any threshold above
\(0.55\); the counterparty network displays a closed cycle of six red
focal nodes connected by thick edges; the per-trader LLM cards each
contain the phrase ``ring signature'' or ``cycle structure''. The
investigator's verification is then a question of (i) noting the
cycle in the counterparty network, (ii) confirming the burst-window
shading on the price chart aligns with the LLM card, and (iii)
clicking through to the underlying trades. End-to-end, this is the
workflow the pipeline was designed for.

\subsection{Per-Run Variance on the ABIDES Cohort}
\label{sec:abides-var}
On the expanded ABIDES cohort, the per-run breakdown deserves
attention because three of the fifteen runs are responsible for most
of the spread in clique recall. Across the five ring runs the
augmented network's recall is exactly \(1.000\) with zero variance;
across the five mixed runs it is \(0.98\pm0.02\); across the five
clique runs it is \(0.91\pm0.11\) with one outlier at \(0.69\). The
outlier corresponds to a clique scenario whose front counterparties
happened to receive unusually balanced order flow, partially masking
the burst-concentration signature. This bimodality is a fair argument
for surfacing per-run distributions in the dashboard alongside the
cohort-mean numbers.

\section{Deployment Considerations}
\label{sec:deploy}

The empirical results imply a clear deployment narrative.
\emph{Ranking} generalises across distributions: the augmented network
holds AUC \(0.842\) on ABIDES versus the baseline's chance, and the
bolt-on stack holds AUC \(0.91\)--\(0.99\) leave-one-family-out.
\emph{Operating points} do not generalise: purity collapses in the
cross-generator regime, and the locked threshold tuned on one cohort
flags too many benigns on a structurally different one. The right
production hygiene is therefore per-generator (or per-deployment)
threshold calibration, not per-generator retraining.

The pipeline supports this directly. Threshold calibration is the
cheapest layer to re-run---no model retraining, no graph rebuild, just
a threshold sweep over pre-computed per-trader scores. We argue that
this should be part of any production rollout plan: deploy the
detector with a default threshold drawn from the training cohort, hold
out a small calibration cohort from the operational distribution, and
sweep the threshold on the calibration cohort to find the operating
point that matches the surveillance team's purity floor.

The architectural contrast between the bolt-on and the end-to-end
detector also has a practical consequence. For deployments where the
operational distribution is expected to drift across structural axes
(different exchanges, different ticker classes, different
manipulation taxonomies than were anticipated at training), the
bolt-on stack is the more robust choice. For deployments where the
operational distribution is well-characterised at training and the
goal is to maximise within-distribution recall at high purity, the
augmented network wins comfortably. The pipeline supports both, and
the choice is a deployment decision rather than a methodological one.


\section{Discussion: Why Bolt-On Beats End-to-End on Family Shift}
\label{sec:disc}

The most provocative empirical result of this work is the
leave-one-family-out contrast: the bolt-on classifier retains AUC
\(0.91\)--\(0.99\) on a held-out family, while the augmented network
collapses to chance. We argue this contrast is informative for system
design and is not just a methodological accident.

\paragraph{Trader-marginal vs.\ relational scoring.} The bolt-on
classifier consumes a seven-dimensional per-trader vector and decides
each trader's score independently of the others. Its decision surface
therefore depends only on whether the engineered features describe a
manipulator-like trader \emph{in isolation}. Because the features were
chosen to be taxonomy-grounded but family-agnostic
(\(\phi_1\)--\(\phi_6\) describe what manipulation looks like, not what
clique manipulation looks like), the decision surface transfers cleanly
across families.

\paragraph{Family-specific neighbourhood priors.} The augmented network,
in contrast, performs two layers of message passing before emitting an
edge probability. The message-passing step pools each node's neighbour
feature vectors, which means the network can---and empirically does---learn
that ``a node whose neighbourhood resembles a ring''
or ``a node whose neighbourhood resembles a clique'' is suspicious.
These neighbourhood signatures are family-specific by construction. When
a family is held out from training, the network's structural prior for
that family is absent, and the trader-marginal signal it does see (via
the same six engineered features) is over-ridden by the inductive bias
of the message-passing layers.

\paragraph{Implication for system design.} For deployments where the
operational distribution is expected to drift across structural axes
(a new exchange whose manipulators differ taxonomically from the
training distribution; a regulator-imposed microstructure change that
shifts neighbourhood patterns), the bolt-on stack is the more robust
choice. For deployments where the operational distribution is
well-characterised at training and the goal is to maximise
within-distribution recall at high purity, the augmented network wins
comfortably. The pipeline supports both, and we view the architectural
choice as a deployment-time decision rather than a model-evaluation
decision.

\paragraph{When the contrast might not hold.} Two caveats are
important. First, the v4 family-disjoint result was obtained under the
five-day training schedule that was tuned for the full cohort, and the
schedule's interaction with a reduced cohort is partially responsible
for the v4 collapse. With a longer training schedule and a larger
validation pool (e.g.\ \(20\%\) of per-day runs rather than \(10\%\),
which the reduced cohort already constrains to about three graphs), we
expect v4's family-disjoint performance to recover at least partially.
This is a clear follow-up experiment. Second, the message-passing
inductive bias that hurts the v4 family-disjoint result is the
\emph{same} bias that lifts v4 to AUC \(0.984\) within-distribution and
\(0.842\) on the cross-generator cohort; the two are properties of the
same architecture, not separable knobs.

\paragraph{Operational workflow.} A pragmatic deployment that wants
both robustness and within-distribution lift would run the bolt-on
classifier as the primary detector for flag generation, and use the
augmented network as a confirmation layer in the dashboard. Flags
where both detectors agree could be promoted to top-of-queue;
disagreements could be triaged manually or downweighted. The pipeline
supports this directly since both detectors emit per-trader scores
into the same artifact directory and the dashboard already reads from
there.

\section{Limitations and Future Work}
\label{sec:lim}

The most important limitation of the work is that both substrates are
synthetic. Although the primary synthesiser is calibrated to NSE
bhavcopy first moments and ABIDES adds genuinely different
microstructure, neither substrate captures the full heavy-tailed
event structure of a real market under stress. A labelled real-data
pilot, in which the detector is applied to a slice of historical NSE
data with regulator-confirmed manipulation cases, would be the
strongest possible validation; the absence of such a labelled cohort is
a recurring limitation in the literature, and we do not solve it here.

Three further limitations apply. First, the family-disjoint v4
retrainings used the training schedule tuned for the full cohort, and
the schedule's interaction with a reduced training set and reduced
validation pool is responsible for at least some of the observed
collapse. A separate study with a longer schedule and larger
validation pool is indicated. Second, the LLM justification card is
not formally validated for factual consistency with the underlying
feature values; we constrain the prompt to the supplied feature
context, but a separate evaluation of the justification card's
faithfulness is needed. Third, the operating-point selection
methodology assumes the investigator team has a single purity floor,
which may be an oversimplification of the per-flag triage workflow.


\section{Threats to Validity}
\label{sec:threats}
Four threats to the generalisability of the empirical results deserve
explicit listing. First, the locked decision threshold used in the
within-generator and ABIDES tables is tuned in-sample on the
evaluation cohort, which inflates absolute purity numbers in both
detectors equally; the cross-detector comparison is fair, but the
absolute operating point is mildly optimistic. Second, both substrates
are synthetic. ABIDES adds genuinely different microstructure but is
still a simulator. Real-market manipulators may produce patterns that
the engineered features were not designed for. Third, the
family-disjoint v4 retraining used the training schedule tuned for the
full five-day cohort, and the schedule's interaction with a reduced
cohort and reduced validation pool is responsible for at least some of
the observed v4 collapse. A retraining with looser early-stopping and a
larger validation pool would be a more definitive test of the
architectural contrast. Fourth, the LLM justification card has been
constrained by prompt structure rather than validated for factual
consistency with the underlying feature values; an independent
faithfulness evaluation is required before the card is treated as
audit evidence.

\section{Reproducibility}
\label{sec:repro}
All artifacts---calibration manifests, synthetic cohorts, feature
caches, model checkpoints, dashboard configurations---are persisted as
plain files in a single \texttt{outputs/} directory, versioned by
cohort name and detector variant. The training pipeline runs entirely
inside Docker containers; a single command (\texttt{docker compose run})
launches the trainer container against the workstation GPU. The bolt-on
classifier runs on CPU outside Docker as a standalone Python script and
completes in under a minute. The dashboard is one Streamlit process
pointed at the same \texttt{outputs/} directory.

\section{Conclusion}
\label{sec:concl}

We have presented a surveillance pipeline for trade-market
manipulation detection that ties together five components: a daily
NSE calibration loop, a calibrated synthetic substrate, a second
substrate (ABIDES) used for cross-generator stress testing, a
modular detector layer comprising a graph neural network with optional
second-stage classifier, and an investigator dashboard that surfaces
flagged traders with feature-level evidence and a language-model
justification card. The pipeline lifts the detector's within-generator
out-of-distribution edge AUC from \(0.638\) to \(0.984\), achieves
pooled trader AUC \(0.968\) in the bolt-on configuration, and retains
edge AUC \(0.842\) on a cross-generator cohort drawn from a different
simulator. A leave-one-family-out evaluation reveals an unexpected
architectural contrast: the bolt-on stack maintains AUC
\(0.91\)--\(0.99\) on the held-out family while the augmented network
collapses to chance, suggesting that for deployments where structural
distribution shift is expected, the simpler classifier is a more
robust choice. The pipeline supports per-generator threshold
calibration as the recommended deployment hygiene for the
cross-distribution operating-point issue. We argue the pipeline is the
contribution: each layer was motivated by a specific surveillance
requirement, and the interfaces between layers are simple enough that
individual components can be swapped without disturbing the rest.

\begin{thebibliography}{99}

% Post-audit bibliography (v2, 2026-08-06). Ten entries; all high-confidence
% real and independently DOI-verifiable. Ten problematic entries removed
% during Workstream A of the pre-submission audit remediation:
%   khanna2021review  — journal "Journal of Financial Surveillance" not
%                        verifiable in any Springer/ACM/IEEE index
%   sun2022fraudgnn   — key/content correspondence unverifiable
%   cong2020amlsim    — attribution wrong (AMLSim is Suzumura & Kanezashi)
%   nguyen2023iclr    — cited ICLR workshop "ML for Finance" does not exist
%   leskovec2014snap  — key year 2014 vs entry year 2016 mismatch
%   chen2016xgboost, wang2019fdgars, liu2021pickfeat, xu2019gin,
%   dou2020carefraud — uncited in v1; will be re-added in Workstream E1
%                       when concrete external-baseline comparisons land
% See docs/AUDIT.md § B10-B11 and docs/REMEDIATION_PLAN.md workstream A.

\bibitem{cao2014pso}
Y. Cao, Y. Li, S. Coleman, A. Belatreche, and T. McGinnity,
``Particle-swarm-optimised support-vector machine for detecting
market manipulation,'' in \emph{Proc. IEEE Symp. on Computational
Intelligence for Financial Engineering}, 2014.

\bibitem{hamilton2017graphsage}
W. L. Hamilton, R. Ying, and J. Leskovec, ``Inductive representation
learning on large graphs,'' in \emph{Advances in Neural Information
Processing Systems (NeurIPS)}, 2017.

\bibitem{velickovic2018gat}
P. Veli\v{c}kovi\'{c}, G. Cucurull, A. Casanova, A. Romero, P. Li\`{o},
and Y. Bengio, ``Graph attention networks,'' in \emph{Int. Conf. on
Learning Representations (ICLR)}, 2018.

\bibitem{kipf2017gcn}
T. N. Kipf and M. Welling, ``Semi-supervised classification with graph
convolutional networks,'' in \emph{Int. Conf. on Learning
Representations (ICLR)}, 2017.

\bibitem{lin2017focal}
T.-Y. Lin, P. Goyal, R. Girshick, K. He, and P. Doll\'{a}r,
``Focal loss for dense object detection,'' in \emph{Proc. IEEE Int.
Conf. on Computer Vision (ICCV)}, 2017.

\bibitem{friedman2001gradient}
J. H. Friedman, ``Greedy function approximation: a gradient boosting
machine,'' \emph{Annals of Statistics}, vol. 29, no. 5, pp. 1189--1232,
2001.

\bibitem{recht2019imagenet}
B. Recht, R. Roelofs, L. Schmidt, and V. Shankar, ``Do ImageNet
classifiers generalize to ImageNet?'' in \emph{Proc. Int. Conf. on
Machine Learning (ICML)}, 2019.

\bibitem{koh2021wilds}
P. W. Koh, S. Sagawa, H. Marklund, S. M. Xie, M. Zhang, A. Balsubramani,
W. Hu, M. Yasunaga, R. L. Phillips, I. Gao, T. Lee, E. David, I. Stavness,
W. Guo, B. A. Earnshaw, I. S. Haque, S. Beery, J. Leskovec, A. Kundaje,
E. Pierson, S. Wang, D. Koller, S. Levine, and P. Liang, ``WILDS: A
benchmark of in-the-wild distribution shifts,'' in \emph{Proc. Int.
Conf. on Machine Learning (ICML)}, 2021.

\bibitem{byrd2019abides}
D. Byrd, M. Hybinette, and T. H. Balch, ``ABIDES: Towards
high-fidelity multi-agent market simulation,'' in \emph{Proc. ACM
SIGSIM Conf. on Principles of Advanced Discrete Simulation}, 2019.

\bibitem{cont2001stylized}
R. Cont, ``Empirical properties of asset returns: stylized facts and
statistical implications,'' \emph{Quantitative Finance}, vol. 1,
no. 2, pp. 223--236, 2001.



\bibitem{streamlit2024}
Streamlit Inc., ``Streamlit: A faster way to build and share data apps,''
Software, 2024. \url{https://streamlit.io/}

\bibitem{ollama2024}
Ollama, ``Ollama: get up and running with large language models locally,''
Software, 2024. \url{https://ollama.com/}

\bibitem{docker2024}
Docker Inc., ``Docker: Accelerated container application development,''
Software, 2024. \url{https://www.docker.com/}

\end{thebibliography}

\end{document}
